\NativeSlot{45}{61}{%
\ExamGroup{二、选择（将正确答案的序号填在括号里）。{\normalfont\kaishu\fontsize{10.8bp}{14bp}\selectfont （每题2分，共12分）}}

% p026_s01

\ExamQuestion{1.}{5}%
计算\(32\times125\)，下面\mbox{（\hspace{1.8em}）}中计算方法是错误的。\par
\ExamOptions{A. \(125\times8+125\times4\)}{B. \(125\times8\times4\)}{C. \(125\times30+125\times2\)}{D. \(125\times40-125\times8\)}
}
\NativeSlot{148}{43}{%
% p026_s02

\ExamQuestion{2.}{5}%
下面的算式中，正确的是\mbox{（\hspace{1.8em}）}。\par
\ExamOptions{A. \(54\times13+7=54\times(13+7)\)}{B. \(99\times16=100\times16-1\)}{C. \(350\div2\div5=350\div(2\times5)\)}{D. \(45\times36=63\times54\)}
}
\NativeSlot{225}{117}{%
% p026_s03

\ExamQuestion{3.}{5}%
算式“\(783-256-\square\)”可以用减法的性质进行简算，则方框表示的数是\mbox{（\hspace{1.8em}）}。\par
\ExamOptions{A. 117}{B. 256}{C. 283}{D. 244}

% p026_s04

\ExamQuestion{4.}{5}%
以下选项\mbox{（\hspace{1.8em}）}的画图方式，可以说明\((2+3)\times m=2\times m+3\times m\)。\par
\ExamOptions{A. 线段图}{B. 分割长方形面积图}{C. 文具总价图}{D. 两个长方形宽均m、高5的图}
\BookFigure{figures/p026_s04-01.tex}
}
\NativeSlot{423}{73}{%
% p026_s05

\ExamQuestion{5.}{5}%
在古埃及，人们采用“双倍法”计算乘法。如计算\(23\times14\)的过程：\(23\times2=46\)，\allowbreak{}\(23\times4=92\)，\allowbreak{}\(23\times8=184\)，\allowbreak{}\(23\times14=46+92+184=322\)。实际上“双倍法”计算使用的运算律是\mbox{（\hspace{6em}）}。\par
\ExamOptions{A. 乘法结合律}{B. 乘法交换律}{C. 乘法分配律}{D. 乘法结合律和分配律}
}
\NativeSlot{548}{75}{%
% p026_s06

\ExamQuestion{6.}{5}%
如右图，要计算这个图形的面积，下面算式中正确的有\mbox{（\hspace{1.8em}）}个。①\(15\times32+15\times18\)；②\(15\times32+18\)；③\(15\times(32+18)\)；④\((15+18)\times32-18\times(32-15)\)。\par
\ExamOptions{A. 1}{B. 2}{C. 3}{D. 4}
\BookFigure{figures/p026_s06-01.tex}
}
\NativeSlot{674}{70}{%
\ExamGroup{三、计算。{\normalfont\kaishu\fontsize{10.8bp}{14bp}\selectfont （共35分）}}
\noindent 1.\quad 直接写得数。\par
\noindent%
% p026_s07
\NativeOralCell{45+55=}%
% p026_s08
\NativeOralCell{125\times8=}%
% p026_s09
\NativeOralCell{5\times7\times4=}%
% p026_s10
\NativeOralCell{25\times6\times4=}%
\par
\noindent%
% p026_s11
\NativeOralCell{4\times25=}%
% p026_s12
\NativeOralCell{123-97=}%
% p026_s13
\NativeOralCell{6\times4\div4\times6=}%
% p026_s14
\NativeOralCell{130-30\div5=}%
\par

}
